\section{Related Work}
\label{sec:related_work}

Our work is positioned at the intersection of parameter-space model merging, dynamic routing/Mixture of Experts (MoE), and statistical regularization under low-data calibration constraints. 

\subsection{Parameter-Space Model Merging}
Parameter-space model merging has emerged as a prominent research direction to consolidate multiple task-specific expert networks into a single model without the computational cost of multi-task joint pretraining. The foundation of this line of research lies in Task Vectors \cite{ilharco2022editing}, which demonstrated that the difference between a fine-tuned model and its pre-trained base model can be treated as a coherent task-specific parameter direction. By scaling and adding these vectors to the base model, multiple capabilities can be blended. 

Subsequent work addressed the destructive interference that occurs when merging overlapping parameters. TIES-Merging \cite{yadav2023ties} introduced sign consensus, parameter trimming, and scaling to resolve parameter conflicts. DARE \cite{yu2024dare} utilized random pruning and scaling to merge models with up to 99\% parameter reduction. ZipIt \cite{stoica2023zip} matched and merged correlated features across models, while Fisher-weighted merging \cite{matena2022merging} used diagonal Fisher information matrices to guide parameter weighting. Despite these advances, these classical merging methods are strictly \emph{static}—they apply a single, global, fixed coefficient per model, which cannot adjust to sample-by-sample differences during inference. Moreover, as we formalize in Appendix~\ref{app:nonlinear}, even if these non-linear alignment methods are extended to dynamic, test-time routing settings, they remain bound to the same fundamental Dynamic Routing Paradox due to covariance estimation singular-value collapse in data-scarce splits.

More recently, training-free dynamic merging methods, such as DAWIN \cite{dawin2024} and SE-Merging \cite{semerging2025}, have been proposed to adapt merging coefficients on-the-fly during inference without requiring a separate, offline calibration phase. These methods typically compute routing weights dynamically by analyzing test-batch statistics or activation distributions on the current input stream. However, we note that these training-free approaches are heavily reliant on test-batch properties. When deployed in true, sample-wise vectorized pipelines (where the batch size is $B=1$ for low-latency interactive streaming), the absence of a batch population prevents these methods from computing robust test-batch statistics. As a result, training-free dynamic merging is highly vulnerable to severe performance degradation or high latency overhead when operating on individual, single-sample streams ($B=1$). In contrast, our Prior-Driven Classical Routing Framework addresses this bottleneck by utilizing a compact, offline-calibrated router with a robust maximum-entropy uniform prior, which guarantees high stability and optimal latency regardless of the inference batch size.

\subsection{Dynamic Routing and Mixture of Experts}
Dynamic routing and Mixture of Experts (MoE) architectures allow models to dynamically route inputs to specialized sub-networks or "experts" based on input-dependent routing networks \cite{shazeer2017outrageously}. Sparse MoEs, such as GShard \cite{lepikhin2020gshard} and Switch Transformers \cite{fedus2022switch}, activate a subset of expert layers to scale the model's capacity while keeping the active compute constant. 

In the context of model merging, Layer-wise Low-dimensional Classical Routing (L3 routing) \cite{muqeeth2023l3} introduced dynamic routing directly to parameter merging, enabling input-dependent activation of task vectors on a layer-by-layer basis. Our work builds upon the L3 routing paradigm but exposes a fundamental confounder: \emph{heterogeneity collapse}. Under mixed-task streaming inputs, standard batch-averaging of dynamic routing coefficients destroys the sample-specific routing benefits, collapsing the weights back to uniform compromise coefficients. We resolve this through task-variance regularization and batch-independent vectorized routing.

\subsection{Quantum-Inspired Machine Learning and Wave Superposition}
To capture highly non-linear task relationships and resolve complex optimization paths in model merging, recent works have turned to quantum-inspired mathematical models. The state-of-the-art Quantum Waveface Superposition (QWS-Merge) \cite{qwsmerge2025} employs complex wave phase equations and wave cosine phase-interferometry to compute merging weights. 

While quantum-inspired machine learning (QML) has shown theoretical promise in modeling high-dimensional correlations \cite{biamonte2017quantum}, our empirical evaluation reveals that wave-interference equations introduce highly non-monotonic activation landscapes. Under severe data scarcity (e.g., 64 calibration samples), this complex formulation is highly sensitive, leading to a performance drop under vectorized sample-wise streaming ($56.19\%$ under $B=1$ compared to the static Uniform baseline of $58.00\%$). Our work advocates for proper classical statistical regularization and prior-layer design, which we empirically demonstrate provides superior and stable generalization under vectorized deployment.
